\section{Static rigidity and the original horizon embedding}
\label{sec:static}

The preceding variational argument has located the area multiplier on
$S$, the original smooth boundary.  We now turn that local information
into a global description of the original metric and second fundamental
form.  This step requires more than solving a vacuum static equation:
the adjoint initially permits null dust, and the norm of its stationary
field may vanish away from $S$.

We use the strong-decay equality class of
Definition~\ref{def:setting:rigidity}: $n\geq3$, the exterior is complete
with its connected boundary included, its complement to a single
Euclidean end is compact, and the decay is $O_6(r^{-q})/O_5(r^{-1-q})$
with $(n-2)/2<q<n-2$.  The dominant energy condition and the
future-normal convention remain those of Definition~\ref{def:data}.
Theorem~\ref{thm:variation:boundary-support} supplies the localized
multiplier; Proposition~\ref{prop:variation:multiplier} supplies its
charge normalization, its complementary support, and its variational
identity for compact variations up to $S$.  In particular, no frontier
carrying the area multiplier remains in the open exterior.  Put
\begin{equation}\label{eq:static:constants}
 r_h=(A/\omega)^{1/k},\qquad m=\tfrac12r_h^{k-1},\qquad
 b^0=E/m,\quad b=-P/m,\qquad c=(k-1)/r_h,\quad\kappa=c/2.
\end{equation}
Thus $(b^0)^2-|b|^2=1$ and $\kappa>0$.  We retain the convention that
$\nu$ points from $S$ into the exterior.

\begin{theorem}[Recovery of the original data]\label{thm:static:classification}
Under the equality and rigidity hypotheses of Theorem~\ref{thm:rigidity},
the original exterior, with its boundary included, admits a smooth
spacelike embedding into the regular Schwarzschild--Tangherlini exterior
of mass $m$.  Its induced metric and future-normal second fundamental
form are $g$ and $K$.  The boundary is a full section of the future
horizon, or the bifurcation sphere, and the end approaches spatial
infinity.
\end{theorem}

Here is the order of the argument.  The adjoint gives smooth fields
$u,X$ and a stationary spacetime containing the original slice.  A twist
identity, tested through the possible zero set of $u^2-|X|^2$, makes
that spacetime locally static wherever the stationary field is timelike.
Write $h$ and $\lambda$ for the resulting static base metric and lapse.
Two auxiliary completions then have different purposes.  After the twist vanishes, the circle metric
$h+\lambda^2dT^2$ rules out an interior zero frontier of the Killing
norm by completeness, Ricci flatness, and the splitting theorem.
We then double the static base conformally across its attached
boundary, compactify one end, and prove that the resulting zero-mass
metric is Euclidean.  This identifies $(h,\lambda)$; a global time
primitive finally recovers the original pair $(g,K)$, including its
smooth attachment to the future horizon or bifurcation sphere.
The nonnegative-mass input is proved below from the numerical
Riemannian Penrose theorem and independent smooth-obstacle geometry.
Its conformal correction and all regularity arguments are part of the
present proof.

\subsection{From the localized multiplier to smooth normalized fields}

Our first task is to turn the multiplier into ordinary smooth fields,
including at the boundary, and to identify their constants at infinity.
The finite first-jet system performs both tasks; the causal inequality
rules out the linear growth that this system would otherwise allow.

Write $u,X$ for the scalar and vector moments in
Proposition~\ref{prop:variation:multiplier}, using the original volume
density.  Their distributional inequalities are
\begin{equation}\label{eq:static:complementarity}
 u\geq |X|_g,\qquad u\mu+J(X)=0.
\end{equation}
After Theorem~\ref{thm:variation:boundary-support}, their interior
equations are
\begin{equation}\label{eq:static:adjoint}
 \sym\nabla X=-uK,\qquad
 \Hess u=-\mathcal S+\frac{\tr_g\mathcal S}{k}g,
\end{equation}
where the smooth-coefficient linear expression is
\begin{align}\label{eq:static:S}
 \mathcal S={}&-u\Ric_g+2u(K^2-\tau K)+\mathcal L_XK
 -(\divg X)K-\divg(K(X,\cdot))g-\sym(J\otimes X^\flat).
\end{align}
Here $K^2_{ij}=K_i{}^aK_{aj}$, and the last divergence is the
divergence of the one-form $K(X,\cdot)$.  These formulas also specify
all the zeroth-order terms needed in the continuation argument.

\begin{lemma}\label{lem:static:smooth-fields}
The interior moments are smooth and extend smoothly to $S$.  The
multiplier has no additional scalar or vector measure supported on
$S$.
\end{lemma}

\begin{proof}
The first equation in Equation~\eqref{eq:static:adjoint} has the usual
finite prolongation.  More explicitly, if
$T_{ij}=\nabla_iX_j+\nabla_jX_i=-2uK_{ij}$, the combination
\[
 \nabla_iT_{j\ell}+\nabla_jT_{i\ell}-\nabla_\ell T_{ij}
\]
equals $2\nabla_i\nabla_jX_\ell$ plus curvature contractions with $X$;
this follows by commuting the two derivatives in the three terms.
Consequently $\nabla^2X$ is a smooth linear combination of
$u,X,du$.  Together with the second equation in
Equation~\eqref{eq:static:adjoint}, this is a closed first-order
homogeneous system for
\[
 \mathcal Y=(u,X,du,\nabla X),\qquad
 \nabla_i\mathcal Y=\mathcal A_i\mathcal Y,
\]
with smooth coefficients.  The interior regularity already obtained
for the moments in the preceding section, followed by these equations,
gives smoothness.  Alternatively their traced equations and smooth
elliptic regularity give the first bootstrap directly.

Fix an interior slice of a compact collar of $S$.  Along each collar
normal, the displayed system is an ordinary linear differential
equation with smooth coefficients on a closed finite interval.  Its
solution and all derivatives with respect to the footpoint extend to
the endpoint.  The original solution agrees with this extension by
uniqueness along each normal.  The tangential equations continue to
hold there by continuity.  This proves smooth extension without yet
discarding measures on the boundary.

Now use metric variations $h$ whose value and first derivatives vanish
on $S$, with $p=0$.  Their area and expansion variations vanish there.
For the smooth interior fields, integration by parts has zero boundary
term and zero interior adjoint term.  On the other hand the scalar
curvature variation at $S$ can be prescribed freely: in Gaussian
coordinates take $h_{AB}=s^2\psi(y)g_{AB}|_{s=0}$ times a collar
cutoff.  Its double-divergence minus Laplacian-of-trace at $s=0$ is
$-2k\psi$, whereas all first-jet and algebraic terms vanish.  The
multiplier identity therefore annihilates every smooth scalar test
against the scalar boundary measure.  That measure is zero.  The
vector boundary measure is dominated by it by
Equation~\eqref{eq:static:complementarity}, and is zero as well.
\end{proof}

The following elementary asymptotic lemma also applies to Killing
lapse and shift on the hypersurfaces in the converse.

\begin{lemma}[Causal first jets on an end]\label{lem:static:end-jets}
Suppose $g-\delta=O_2(r^{-q})$, $q>1/2$, and a smooth pair
$Y=(u,X)$ on the coordinate end satisfies $u\geq|X|_g$ and
\begin{equation}\label{eq:static:jet-inequality}
 |D^2Y|\leq C r^{-1-q}|DY|+C r^{-2-q}|Y|.
\end{equation}
Then $Y=Y_\infty+O_2(r^{-q_0})$ for every $0<q_0<q$.
\end{lemma}

\begin{proof}
All estimates are uniform in the angular variable.  Along radial rays,
$|Y(r)|\leq C+\int_{R_0}^r|DY(t)|\,dt$.  Integrating
Equation~\eqref{eq:static:jet-inequality}, taking the supremum on the
sphere, and interchanging the two integrals gives
\[
 \sup_{R_0\leq t\leq r,\,\vartheta}|DY(t,\vartheta)|
 \leq C+C\int_{R_0}^r t^{-1-q}
       \sup_{R_0\leq s\leq t,\,\vartheta}|DY(s,\vartheta)|\,dt.
\]
The term containing $Y$ has the required bound because
\[
 \int_s^r t^{-2-q}\,dt\leq C s^{-1-q}.
\]
Gronwall's inequality gives $DY=O(1)$ and $Y=O(r)$, hence
$D^2Y=O(r^{-1-q})$.  Each Cartesian derivative has a limit along
every ray.  Comparing two rays on a sphere of radius $r$, along a
spherical path of length at most $C r$, shows that their derivative
values differ by $O(r^{-q})$.  Thus all ray limits form a single
constant matrix $B$, and $Y(x)=Bx+o(r)$ uniformly on spheres.

The first row of $B$ is zero: its linear function is nonnegative on
the entire unit sphere, and replacing a direction by its negative
changes its sign.  Dividing $u\geq|X|_g$ by $r$ now forces every
remaining row of $B$ to vanish.  Terminal integration of $D^2Y$ gives
$DY=O(r^{-q})$.

If $q>1$ this already makes $Y$ bounded.  If $1/2<q<1$, radial
integration first gives $Y=O(r^{1-q})$, and reinsertion gives
$D^2Y=O(r^{-1-2q})$ and, since the derivative limit is zero,
$DY=O(r^{-2q})$.  This is integrable because $2q>1$, so $Y$ is
bounded.  If $q=1$, replace the exponent in this intermediate step
by any number strictly between $1/2$ and $1$.
In all cases choose an intermediate $\widehat q$ with
$1/2<\widehat q<q$ (or use $q$ itself away from $1$) so that
$DY=O(r^{-a})$ for some $a>1$.  Boundedness of $Y$ in
Equation~\eqref{eq:static:jet-inequality} then gives
$D^2Y=O(r^{-2-q})$, followed by $DY=O(r^{-1-q})$ and convergence
of $Y$ on each ray with error $O(r^{-q})$.  Angular comparison using
this last derivative bound shows that the limits of $Y$ are the
same.  Weakening the exponent proves the stated $O_2$ assertion.
\end{proof}

The prolonged system satisfies Equation~\eqref{eq:static:jet-inequality}:
curvature, $\nabla K$, $K^2$ and $J$ multiply values, while $K$ and
the coordinate Christoffel symbols multiply first derivatives.
The assumed end differentiability suffices for the two differentiated
estimates.  Fix henceforth
\begin{equation}\label{eq:static:q0}
 (n-2)/2<q_0<q,\qquad q_0\ne1.
\end{equation}
The charge prototypes in Proposition~\ref{prop:variation:multiplier}
identify the constants.  After integrating the adjoint equations by
parts, the only surviving terms at infinity are the constant lapse
times the scalar-curvature flux and twice the constant shift times
the momentum-constraint flux.  For the scalar prototype
$h=O(r^{2-n})$, $Dh=O(r^{1-n})$, terms such as
$u_\infty(g-\delta)Dh$, $u_\infty\Gamma h$, and
$KX_\infty h$ have integrated flux $O(r^{-q})$.
Replacing a constant field by its remainder, or using $DY$ in a
boundary term, contributes $O(r^{-q_0})$.  The vector prototypes
have the same or better orders.  Thus all background and field
errors have total flux $O(r^{-q})+O(r^{-q_0})=o(1)$.
The scalar prototype changes $E$ nontrivially
and the vector prototypes span the changes of $P$.  Comparing their
coefficients in the multiplier identity therefore yields
\begin{equation}\label{eq:static:normalization}
 u=b^0+O_2(r^{-q_0}),\qquad X=b+O_2(r^{-q_0}).
\end{equation}
Finally $u>0$ in the open exterior.  At an interior zero of $u$,
nonnegativity gives $du=0$, and domination gives $X=0$ and $DX=0$.
The full first jet is then zero; uniqueness in the homogeneous
system along piecewise smooth paths contradicts
Equation~\eqref{eq:static:normalization}.

\subsection{The stationary spacetime and its horizon boundary laws}

The fields now have a nonzero future-causal limit and a positive lapse
in the interior.  We use them to construct a stationary spacetime whose
$z=0$ slice has exactly the given $g$ and $K$.  The multiplier identity
will determine both the possible matter tensor and the data at $S$.

On $\R_z\times\operatorname{int}\Omega$ define
\begin{equation}\label{eq:static:stationary-metric}
 \mathbf g=-u^2dz^2+g_{ij}(dx^i+X^i dz)(dx^j+X^j dz),
 \qquad \xi=\partial_z=u\mathbf n+X.
\end{equation}
We orient time by $\mathbf n=u^{-1}(\partial_z-X)$.
Its slice second form is
$-\sym\nabla X/u=K$, with precisely the future-normal convention
of the theorem.

\begin{lemma}\label{lem:static:stress-boundary}
The Einstein tensor and the boundary values satisfy
\begin{align}\label{eq:static:dust}
 \mathbf G&=\mu u^{-2}\xi^\flat\otimes\xi^\flat,
 &\Ric_{\mathbf g}(\xi,\cdot)&=0,\\
 \label{eq:static:boundary-laws}
 X|_S&=u\nu,&\partial_\nu u+K(X,\nu)&=\kappa.
\end{align}
On the connected $S$, either $u>0$ everywhere or $u=0$ everywhere.
Writing $W=u^2-|X|_g^2$, a deleted collar of $S$ has $W>0$.
\end{lemma}

\begin{proof}
In the open exterior, where $u>0$, Equation~\eqref{eq:static:adjoint}
gives
\[
 b=-u^{-1}\sym\nabla X=K,\qquad
 \Hess u-(\Delta u)g+\mathcal S=0.
\]
Apply Lemma~\ref{lem:variation:stationary-action} with this $b$.
Complementarity eliminates its volume term and yields
\[
 u\mathbf G_{ij}=-J_{(i}X_{j)}.
\]
Gauss--Codazzi supplies $\mathbf G(\mathbf n,\mathbf n)=\mu$ and
$\mathbf G(\mathbf n,\partial_i)=J_i$.

If $\mu>0$, equality in
$0=u\mu+J(X)\geq\mu(u-|X|)$ forces
$|X|=u$ and $J=-\mu X^\flat/u$.  If $\mu=0$, DEC gives $J=0$.
Since $\xi^\flat(\mathbf n)=-u$ and its spatial part is
$X^\flat$, these identities give the first part of
Equation~\eqref{eq:static:dust}.  In particular $\mu W=0$.
The stress tensor has zero spacetime trace and annihilates $\xi$;
the same is therefore true of its Ricci contraction.

For compact variations extending to $S$, the multiplier identity
now reads
\[
 2k\omega\mathcal E'-\frac c2\int_S\tr_S h\,dA
 =\int_\Omega\{u(2\mu)'+2J'(X)-J_i h^i{}_jX^j\}\,dV_g
 -\int_S\theta_+'\,d\zeta.
\]
The free $p=K'$ boundary term, with integration normal $-\nu$, is
$2X_\nu\tr_S p-2p(X_{\rm tan},\nu)$.
Since $\theta_+'=\tr_S p$ in this test, arbitrary normal-tangential
and tangential components give
$X_{\rm tan}=0$ and $d\zeta=2X_\nu dA$.
For a compact conformal test $(h,p)=(2s g,sK)$, the constraint
variation on active rays is
$2k(-\Delta s+K(\nabla s,v))$ and
$\theta_+'=k\partial_\nu s$.  Its boundary identity becomes
\[
 -kc\int_Ss\,dA
 =2k\int_S\{u\partial_\nu s-
       (\partial_\nu u+K(X,\nu))s\}\,dA
       -k\int_S\partial_\nu s\,d\zeta.
\]
The independent value and normal derivative of $s$ give
$d\zeta=2u\,dA$ and Equation~\eqref{eq:static:boundary-laws}.

It remains to justify the dichotomy, since $u$ might initially
vanish at only part of $S$.  One-sided area minimization implies
$H\geq0$: test the area by every nonnegative outward speed.
Use Lie derivative variations along a compactly supported vector field
$Y$ equal to $s\nu$ on $S$, extended across its collar.  Fix an interior
active ray $(x,v)$, let $\phi_t$ be the local flow of $Y$, and put
$(g_t,K_t)=(\phi_t^*g,\phi_t^*K)$.  If $v_t$ is the isometrically
identified ray for $g_t$, then $w_t=\mathrm d\phi_t(v_t)$ satisfies
$|w_t|_{g,\phi_t x}=|v|_{g,x}\leq1$.  Naturality gives the constraint
derivative as $2F'(0)$, where
\[
 F(t)=\mu(\phi_t x)+J_{\phi_t x}(w_t)\geq0.
\]
Since $x$ is interior, this smooth function is defined for both signs
of sufficiently small $t$.  Activity gives $F(0)=0$, so $F'(0)=0$.
This argument also applies when $|v|_g=1$, since it differentiates only
along the transported curve in the closed unit-ball bundle.  Boundary
rays contribute no integral, since Lemma~\ref{lem:static:smooth-fields}
has already eliminated the boundary constraint measure.  The area
variation is $\int_S Hs$, and the expansion variation is
$L_+s$, the outward normal linearization of $\theta_+$.
Thus
\begin{equation}\label{eq:static:boundary-adjoint}
 L_+^*u=\frac c2 H\geq0.
\end{equation}
The operator has principal part $-\Delta_S$ and smooth bounded
coefficients.  The strong minimum principle for a nonnegative
supersolution applies locally: replace its zeroth-order coefficient
by a larger nonnegative constant to obtain the usual form, then
propagate along overlapping balls.  Since $S$ is connected,
$u>0$ throughout $S$, or $u$ is identically zero there.

If $u>0$ on $S$, the tangential derivatives of $W$ vanish, and
the normal-normal component of $\sym\nabla X=-uK$ gives
\begin{equation}\label{eq:static:positive-boundary}
 dW|_S=2\kappa X^\flat=2\kappa u\nu^\flat.
\end{equation}
Thus $W$ is a positive defining function in the exterior collar.
If $u=0$ on $S$, then $X=0$, all its tangential derivatives vanish,
and its symmetrized equation successively gives
$\nabla_\nu X=0$ in the normal-normal and normal-tangential
components.  In original exterior distance $s$,
\begin{equation}\label{eq:static:zero-boundary}
 u=s a,\qquad X=s^2Y_*,\qquad a|_S=\kappa,
 \qquad W=s^2(a^2-s^2|Y_*|_g^2),
\end{equation}
with smooth coefficients.  This also proves positivity in a
deleted collar.
\end{proof}

\subsection{A common stationary-field construction of the static base}
\label{subsec:static:common-base}

We isolate the part of the stationary argument that also applies to
the weaker three- and four-dimensional end hypotheses.  Its output
allows complete ends arising from an interior null set.  The later
classification arguments will exclude those ends by different methods.
The twist estimate is the main step: near an arbitrary null set, a
logarithmic bound controls the transverse derivatives in the flux.

\begin{proposition}[The complete attached static base]
\label{prop:static:common-base}
Let $(\Omega,g)$ be a smooth connected $n$-manifold, $n\geq3$,
complete with its nonempty compact smooth boundary
$S=\bigsqcup_{i=1}^{\ell}S_i$, $1\leq\ell<\infty$, included.
Suppose the complement of a compact
set is one Euclidean coordinate end.  Let $K,u,X$ be smooth up to
the one-sided boundary, and suppose that, for some $q_0>(n-2)/2$,
\[
 g-\delta=O_2(r^{-q_0}),\qquad
 (u,X)=(b^0,b)+O_2(r^{-q_0}),\qquad
 (b^0)^2-|b|^2=1,\quad b^0>0.
\]
Assume $u>0$ in the open exterior, $u\geq|X|_g$, and
$\sym\nabla X=-uK$.  On $\R\times\operatorname{int}\Omega$, set
\[
 \mathbf g=-u^2dz^2+g_{ij}(dx^i+X^i dz)(dx^j+X^j dz),
 \qquad \xi=\partial_z,
\]
and assume $\Ric_{\mathbf g}(\xi,\cdot)=0$ throughout this
interior spacetime and
$\Ric_{\mathbf g}=0$ where $W:=u^2-|X|_g^2>0$.
On each connected boundary component $S_i$, suppose
\[
 X=u\nu,\qquad \partial_\nu u+K(X,\nu)=\kappa_i,
 \qquad \kappa_i>0\ \hbox{constant},
\]
where $\nu$ points into $\Omega$, and suppose either $u>0$
throughout $S_i$ or $u=0$ throughout it.

Put $\alpha=X_g^\flat$.  On $\mathcal U=\{W>0\}$ define
\begin{equation}\label{eq:static:orbit-metric}
 \begin{gathered}
 \lambda=\sqrt W,\qquad
 h_s=g+W^{-1}\alpha\otimes\alpha,\qquad
 C=h_s^{-1}=g^{-1}-u^{-2}X\otimes X,\\
 A_s=W^{-1}\alpha,\qquad F_s=dA_s.
 \end{gathered}
\end{equation}
Then $\mathcal U$ is connected and contains the end and a deleted
collar of every $S_i$.  One has $F_s=0$, and, writing $h=h_s$,
\begin{equation}\label{eq:static:common-equations}
 \Hess_h\lambda=\lambda\Ric_h,\qquad
 \Delta_h\lambda=0,\qquad R_h=0,\qquad 0<\lambda<1.
\end{equation}
The metric $h$ has a smooth one-sided attachment at each $S_i$,
with induced metric $g|_{TS_i}$, zero second fundamental form,
and $\partial_{\nu_h}\lambda=\kappa_i$.  The attached base is
complete with its boundary included.  Every possible interior null
frontier in the original exterior is infinitely far away for $h$.
The even reflection of $h$ and odd reflection of $\lambda$ are
$C^{2,\alpha}$ for every $0<\alpha<1$, and the metric double is
complete.  After a fixed linear coordinate change, the given end
of $h$ is asymptotically Euclidean with $O_2(r^{-q_0})$ decay.
\end{proposition}

The proposition does not yet assert $W>0$ on the entire original
exterior or exactness of the closed form $A_s$.

\begin{lemma}[Vanishing of the stationary twist]\label{lem:static:twist}
Under the hypotheses of Proposition~\ref{prop:static:common-base},
$F_s=0$ on every component of $\mathcal U$.
\end{lemma}

\begin{proof}
First the boundary laws and $\sym\nabla X=-uK$ give the collars
needed below.  On a component where $u>0$, differentiating
$W=u^2-|X|^2$ and using its vanishing tangential derivatives gives
\begin{equation}\label{eq:static:common-positive-collar}
 dW|_{S_i}=2\kappa_iX^\flat=2\kappa_i u\nu^\flat.
\end{equation}
On a component where $u=0$, also $X=0$ and its tangential
derivatives vanish.  The normal-normal and normal-tangential
components of the symmetrized equation give $\nabla_\nu X=0$.
In original exterior distance $s$, therefore,
\begin{equation}\label{eq:static:common-zero-collar}
 u=sa,\qquad X=s^2Y_*,\qquad a|_{S_i}=\kappa_i,\qquad
 W=s^2(a^2-s^2|Y_*|_g^2).
\end{equation}
Both alternatives imply $W>0$ in a deleted collar.  Since $W\to1$
at the end, the interior zero set is compact and separated from
all these collars and the distant end.

We first derive an estimate that permits integration through an arbitrary
interior zero set of $W$. The tensor $C=g^{-1}-u^{-2}X\otimes X$ is
smooth and nonnegative throughout the open exterior, including its
zeros of $W$. Let
$D_{\gamma\eta}=\boldsymbol\nabla_\gamma\xi_\eta$; the Killing
equation makes $D$ skew. In the following calculation
$|D|_{\mathbf g}^2$ denotes the full contraction
$D_{\gamma\eta}D^{\gamma\eta}$, which need not be nonnegative.
The Killing wave identity, together with
$\Ric_{\mathbf g}(\xi,\cdot)=0$, gives
$\boldsymbol\square W=-2|D|_{\mathbf g}^2$. The spatial inverse
block of $\mathbf g$ is $C$, its volume density is
$u\sqrt{\det g}$, and $W$ is independent of $z$. Consequently
\begin{equation}\label{eq:static:wave}
 u^{-1}\divg_g(uC\,dW)=-2|D|_{\mathbf g}^2.
\end{equation}
This is an identity of smooth functions through all interior zeros.

Fix a relatively compact neighborhood of such a zero, small enough
that $u$ and $v=|X|_g$ have positive lower bounds there. Write
$e=X/v$, $T=e^\perp$, and choose a local $g$-orthonormal frame
$(e,E_1,\ldots,E_{n-1})$. Let $\mathbf n$ be the future unit normal
to the $z$-slices, so $\xi=u\mathbf n+ve$.
For a spatial index $i$ differentiation of $W=-\mathbf g(\xi,\xi)$
gives
\[
 D_{i\mathbf n}=u^{-1}(-W_i/2-vD_{ie}).
\]
Putting $a_a=D_{ae}$ and $b_{ab}=D_{ab}$, full contraction in this
Lorentz orthonormal frame therefore gives the exact expansion
\[
 -2|D|_{\mathbf g}^2
 =\frac{|dW|_g^2}{u^2}
   +\frac{4v}{u^2}\sum_a W_a a_a
   -\frac{4W}{u^2}\sum_a a_a^2
   -2\sum_{a,b}b_{ab}^2.
\]
On every smaller relatively compact neighborhood, smooth
nonnegativity of $W$ gives $|dW|_g^2\leq C_0W$. This elementary
estimate follows by applying the second-order Taylor bound along
short geodesics in the direction $-\nabla W$. Since $D$ is bounded,
the favorable last two terms imply
\[
 -2|D|_{\mathbf g}^2\leq C_1(W+|d_TW|_g).
\]
Here $d_TW=dW-e(W)e^\flat$. The eigenvalues of $C$ are $1$ on
$T$ and $W/u^2$ along $e$.

Choose a smooth compactly supported cutoff $\eta$ in this
neighborhood. Multiplying Equation~\eqref{eq:static:wave} by
$\eta^2/(W+\delta)$ and integrating gives
\begin{align*}
 E_\delta
 &:={\int}\frac{u\eta^2|dW|_C^2}{(W+\delta)^2}\,dV_g\\
 &=\int\frac{u\eta^2(-2|D|_{\mathbf g}^2)}{W+\delta}\,dV_g
   +2\int\frac{u\eta\langle d\eta,dW\rangle_C}{W+\delta}\,dV_g.
\end{align*}
The preceding upper bound and Young's inequality bound the first
term by $E_\delta/4+C_\eta$: use $W/(W+\delta)\leq1$ and
$|d_TW|_g^2\leq |dW|_C^2$. The same inequality bounds the second
term by $E_\delta/4+C_\eta$, since $C\leq g^{-1}$.
Thus $E_\delta\leq 2C_\eta$, independently of $\delta>0$.
Fatou's lemma gives, wherever $\eta=1$,
\begin{equation}\label{eq:static:log-bound}
 \int_{\{W>0\}} |d_T\log W|_g^2\,dV_g<\infty.
\end{equation}
The integral is local in the fixed neighborhood, including its zero
set as a limiting set. No ellipticity in the direction $e$ at
$W=0$ was required.

The estimate controls the transverse logarithmic derivative without
requiring ellipticity along $X$.  We now derive the twist flux and verify
that it contains only this controlled derivative.  On $\mathcal U$ write $\theta=dz-A_s$, so that
$\mathbf g=-W\theta^2+h_s$ and
\[
 \xi^\flat=-W\theta,\qquad
 d\xi^\flat=-dW\wedge\theta+WF_s.
\]
Raising the index of $\theta$ with $\mathbf g$ gives no spatial
component. Thus the raised spatial block of $d\xi^\flat$ is
$WC^{ik}C^{jl}F_{s,kl}$. The Killing Maxwell identity
$\boldsymbol\nabla_\gamma(d\xi^\flat)^{\gamma\eta}=0$
and stationarity yield
\begin{equation}\label{eq:static:Maxwell}
 \nabla_iQ^{ij}=0,\qquad
 Q^{ij}=uW C^{ik}C^{jl}F_{s,kl},
\end{equation}
where $\nabla$ is the connection of $g$.
Since $C\alpha=(W/u^2)X$, multiplication by $A_s$ and lowering
with $g$ gives
\begin{equation}\label{eq:static:twist-flux}
 \begin{split}
 g_{i\ell}Q^{\ell j}A_{s,j}
 &=\frac1u\left[
 X^j(d\alpha)_{ij}-\frac{|X|_g^2}{W}W_i
       +\frac{X(W)}{W}\alpha_i\right]\\
 &=\frac1u\left[X^j(d\alpha)_{ij}
           -|X|_g^2(d_T\log W)_i\right].
 \end{split}
\end{equation}
The first line is valid even at zeros of $X$ in $\mathcal U$;
the second is used only where $X\ne0$. In particular the flux
$Q^{ij}A_{s,j}$ is exactly perpendicular to $X$.

To convert the divergence identity into vanishing twist, we must
remove three boundaries of integration: the possible interior zero
set, the original horizon, and spatial infinity.  We check their fluxes
separately before taking any limit.  The componentwise collars
\eqref{eq:static:common-positive-collar}--\eqref{eq:static:common-zero-collar}
already separate the compact interior zero set from $S$ and the
distant end.  The boundary estimates below are made on each of
the finitely many components.

Choose a smooth function $f$ which vanishes off a distant end
neighborhood and equals $A_{\infty,i}x^i$ sufficiently far out,
where $A_\infty$ is the constant limit of $A_s$. Put
$\beta=A_s-df$, so $d\beta=F_s$. Near the zero set and $S$,
$\beta=A_s$.

For the interior cutoff take a smooth increasing function
$\chi$, zero on $(-\infty,1]$ and one on $[2,\infty)$, and use
$\chi(W/\delta)$ on a fixed neighborhood of the compact zero set.
Extend it as one outside a slightly larger such neighborhood.
One precise extension is $1-\rho+\rho\chi(W/\delta)$, with
$\rho=1$ near the zero set; for all sufficiently small $\delta$,
the derivative involving $d\rho$ vanishes because $W$ has a
positive lower bound on its support. On
$\{\delta<W<2\delta\}$, transversality in
Equation~\eqref{eq:static:twist-flux} gives
\[
 \left|Q^{ij}A_{s,j}\partial_i\chi(W/\delta)\right|
 \leq C\bigl(|d_T\log W|_g+|d_T\log W|_g^2\bigr)
 \leq C\bigl(1+|d_T\log W|_g^2\bigr).
\]
A finite cover by neighborhoods satisfying
Equation~\eqref{eq:static:log-bound} shows that its integral tends
to zero as $\delta\downarrow0$.

For the boundary cutoff use $\chi(s/\epsilon)$. If $u>0$ on
$S_i$, then uniformly in its collar
\[
 W\asymp s,\qquad e=\nabla s+O(s),\qquad d_TW=O(s).
\]
The flux in Equation~\eqref{eq:static:twist-flux} is bounded and
perpendicular to $e$, so its contraction with $ds$ is $O(s)$.
Since the collar volume is $O(\epsilon)$ and the cutoff derivative
is $O(\epsilon^{-1})$, the boundary error is $O(\epsilon)$.
If $u=0$ on $S_i$, use the first line of
Equation~\eqref{eq:static:twist-flux}. Here
$d\alpha=O(s)$, $d\log W=O(s^{-1})$, and $X=s^2Y$; each term
inside its brackets is $O(s^3)$, whereas $u\asymp s$.
The whole flux is consequently $O(s^2)$, uniformly also where
$X=0$, and the boundary error is $O(\epsilon^2)$.

Finally choose an outer cutoff supported in $\{r<2R\}$ and equal
to one on $\{r<R\}$, with derivative $O(R^{-1})$. On its annulus
the end estimates give
\[
 \begin{gathered}
 \beta=O(r^{-q_0}),\qquad
 Q=O(r^{-1-q_0}),\\
 \left|\int Q^{ij}\beta_j\partial_i\chi_R\,dV_g\right|
 \leq CR^{n-2-2q_0}=o(1).
 \end{gathered}
\]

Let $\zeta$ be the product of these three nonnegative cutoffs.
It has compact support in $\mathcal U$. Testing
Equation~\eqref{eq:static:Maxwell} against $\zeta\beta$ gives
\[
 \int\frac{\zeta uW}{2}
          C^{ik}C^{jl}F_{s,ij}F_{s,kl}\,dV_g
 =-\int Q^{ij}\beta_j\partial_i\zeta\,dV_g.
\]
The left integrand is nonnegative, and strictly positive where
$F_s\ne0$. Fix any compact subset of $\mathcal U$ and choose
the cutoffs to equal one there. First let $\delta\downarrow0$,
then $\epsilon\downarrow0$, and finally $R\to\infty$.
The preceding estimates force the twist energy on that compact
subset to vanish. Exhausting $\mathcal U$ proves $F_s=0$
everywhere on it, without assuming beforehand that the total twist
energy is finite.
\end{proof}

We henceforth write $h=h_s$ and retain the notation $A_s$ and $F_s=dA_s$.


\begin{proof}[Proof of Proposition~\ref{prop:static:common-base}]
We identify the positive component, prove the distance estimate at
its possible interior null frontier, and attach the original boundary.

\smallskip\noindent
\textbf{Step 1: Static equations and the positive component.}
By Lemma~\ref{lem:static:twist}, the one-form
$A_s=X^\flat/W$ is closed on $\{W>0\}$.  On a small simply
connected set write $A_s=d\Phi$ and set $t=z-\Phi$.  Then
\[
 \mathbf g=-\lambda^2dt^2+h.
\]
The vacuum hypothesis applies on this set.  The warped metric
formulas, obtained from
$\mathbf\Gamma^t_{it}=\lambda^{-1}\lambda_i$ and
$\mathbf\Gamma^i_{tt}=\lambda\lambda^i$, are
\[
 \Ric_{\mathbf g,tt}=\lambda\Delta_h\lambda,
 \qquad
 \Ric_{\mathbf g,ij}=\Ric_{h,ij}-\lambda^{-1}\Hess_{h,ij}\lambda.
\]
They give the equations in
Equation~\eqref{eq:static:common-equations} on every positive
component.

There is a positive component $\mathcal C$ containing the end,
since $W\to1$.  Every other positive component has compact
closure in the original exterior: the complement of a coordinate
tail is compact.  Its continuous function $\lambda$ is zero on
its frontier, including any part of $S$.  A positive maximum would
therefore be attained at an interior point, contrary to the strong
maximum principle for its harmonic equation.  Hence there is only
one positive component.  If $\lambda>1$ anywhere in that
component, its limit $1$ at infinity and its zero value on the
original frontier force a maximum at a positive interior point
of a compact original truncation.  The strong maximum principle
excludes this.  Thus $\lambda\leq1$; the boundary collar
prevents constancy, so $0<\lambda<1$.  This argument uses no
regularity of a possible interior frontier.
The positive collar of every $S_i$ belongs to $\mathcal C$.

\smallskip\noindent
\textbf{Step 2: Distance to an interior zero frontier.}
An interior zero of $W$ is infinitely far away for $h$.  Indeed,
on a relatively compact neighborhood of such a zero we have
$u,|X|>0$.  Put $e=X/|X|$ and denote its orthogonal complement
by $T$.  The exact flux identity
Equation~\eqref{eq:static:twist-flux}, with $F_s=dA_s=0$, gives a
uniform bound for $d_T\log W$ on its positive part.  Smooth
nonnegativity gives $|dW|_g^2\leq C W$.  Since
$h^{-1}=g^{-1}-u^{-2}X\otimes X$, it follows that
\begin{equation}\label{eq:static:length-log}
 |d\log W|_{h^{-1}}^2
 =|d_T\log W|_g^2+\frac{(eW)^2}{u^2W}\leq C.
\end{equation}
The integral of $|d\log W|$ along a finite $h$-length curve is
bounded; such a curve cannot have $W\to0$ at an interior point.
Also $h\geq g$.  Every finite-length escape curve has a limit in
the original complete metric space, with $S$ included.  An
interior positive limit is not an escape; an interior zero is
excluded by Equation~\eqref{eq:static:length-log}.  We next
describe precisely the remaining limits on $S$.

\smallskip\noindent
\textbf{Step 3: Attachment at the original boundary.}
Work on a fixed component $S_i$.  In the positive boundary-lapse
case, Equation~\eqref{eq:static:common-positive-collar}
allows original coordinates $(W,y)$ in the collar.  Put
$\alpha=X^\flat$.  Its smooth coefficients satisfy
\[
 \alpha=b(W,y)dW+W\beta_A(W,y)dy^A,
 \qquad b(0,y)=\frac1{2\kappa_i}.
\]
Indeed its tangential coefficients vanish at $W=0$, and the
normal coefficient follows from $dW=2\kappa_i\alpha$ there.
In coordinates $(\lambda,y)$, $W=\lambda^2$, the metric is
\begin{align*}
 h={}&(4\lambda^2g_{WW}+4b^2)d\lambda^2
 +2(2\lambda g_{WA}+2b\lambda\beta_A)d\lambda\,dy^A\\
 &+(g_{AB}+\lambda^2\beta_A\beta_B)dy^A dy^B.
\end{align*}
All coefficient functions on the right are evaluated at
$(\lambda^2,y)$.  Its normal coefficient at zero is
$\kappa_i^{-2}$, its cross coefficients vanish, and its tangential
part is $g|_{TS_i}$.  The $d\lambda^2$ and tangential coefficients
are even smooth functions, and the cross coefficients are odd
smooth functions.  This is a smooth reflected tensor, not merely
continuity of its components.

In the zero boundary-lapse case,
Equation~\eqref{eq:static:common-zero-collar} instead gives directly
\[
 h=g+s^2\frac{Y_*^\flat\otimes Y_*^\flat}
                    {a^2-s^2|Y_*|^2},\qquad
 \lambda=s\sqrt{a^2-s^2|Y_*|^2}.
\]
These are smooth one-sided, the first is positive definite at
$S_i$, and $d\lambda=\kappa_i\,ds$ there.  In both cases the
static equation extends to $S$ and gives $\Hess_h\lambda=0$.
Its tangential restriction is $\kappa_i$ times the second
fundamental form, so that form is zero.  In Gaussian base
distance $\sigma$ the Taylor expansions are therefore
\begin{equation}\label{eq:static:polar-jets}
 h=d\sigma^2+h_0(y)+\sigma^2h_2(y)+O(\sigma^3),\qquad
 \lambda=\kappa_i\sigma+a_3(y)\sigma^3+O(\sigma^4).
\end{equation}
Here and below these one-sided remainders have the corresponding
differentiated Taylor bounds.  The reflected coefficients match
through second derivatives; their third derivatives are locally
bounded on each side.  Thus the even reflection of $h$ and odd
reflection of $\lambda$ are $C^{2,\alpha}$ for every $\alpha<1$.
Limits of finite-length curves on $S$
are limits in this attached collar, as is seen directly from
the coordinates $(\lambda,y)$ or $(s,y)$ just constructed.

The finite-length curve criterion proves completeness of the
attached base, and gluing two copies along its compact boundary
gives a complete metric double.  Finally the original end estimates
give
\[
 h_{ij}=\delta_{ij}+b_i b_j+O_2(r^{-q_0}),\qquad
 \lambda=1+O_2(r^{-q_0}).
\]
A fixed linear change of coordinates takes the constant positive
metric $\delta+b^\flat\otimes b^\flat$ to the Euclidean metric
and preserves these decay orders.
\end{proof}

\subsection{Excluding interior null frontiers by the circle extension}

We return to the all-dimensional equality data.  The common
stationary construction supplies a complete attached static base;
it remains to exclude its possible interior null ends.  Collapsing
a circle at the original horizon gives a complete Ricci-flat
manifold whose product end has incompatible volume growth with a
splitting caused by any additional end.

\begin{proposition}\label{prop:static:positive}
One has $W>0$ throughout $\operatorname{int}\Omega$.  The static
base metric $h=g+W^{-1}X^\flat\otimes X^\flat$ has a smooth
one-sided attachment at $S$, with induced metric $g|_{TS}$, and
\begin{equation}\label{eq:static:static-equations}
 \Hess_h\lambda=\lambda\Ric_h,\qquad
 \Delta_h\lambda=0,\qquad R_h=0,\qquad
 0<\lambda<1,\quad \lambda=\sqrt W.
\end{equation}
At the attached boundary $S$, $|d\lambda|_h=\kappa$ and the
boundary is totally geodesic.  The even reflection of $h$ and odd
reflection of $\lambda$ are at least $C^{2,\alpha}$ for
$0<\alpha<1$.
\end{proposition}

\begin{proof}
The original compact-complement and completeness assumptions,
Equation~\eqref{eq:static:normalization}, and
Lemma~\ref{lem:static:stress-boundary} verify every hypothesis of
Proposition~\ref{prop:static:common-base}.  In particular the
geometric density vanishes where $W>0$, so the stationary metric
is vacuum there.  That proposition gives the connected positive
region $\mathcal C$, its complete attached static base, and all
the asserted local boundary data.  Here the boundary constant is
the single number $\kappa$ in Equation~\eqref{eq:static:constants}.
We need only show that no interior null frontier remains.

\smallskip\noindent
\textbf{Step 1: A smooth complete Ricci-flat circle manifold.}
To rule out missing interior sets, form over $\mathcal C$ the
product with a circle of period $2\pi/\kappa$ and metric
\begin{equation}\label{eq:static:circle}
 \mathfrak h=h+\lambda^2dT^2.
\end{equation}
This product is built from the global static metric and lapse;
it requires no global primitive of $A_s$ or quotient of the
stationary spacetime.
Collapse the circle over $S$ by attaching $S\times D^2$ in
the collar.  This attachment has the required regularity.
To check it explicitly put $\theta=\kappa T$ and
$(z_1,z_2)=(\sigma\cos\theta,\sigma\sin\theta)$.
The normal polar part equals
\[
 dz_1^2+dz_2^2+
 \frac{(\lambda/(\kappa\sigma))^2-1}{\sigma^2}
       (z_1dz_2-z_2dz_1)^2.
\]
By Equation~\eqref{eq:static:polar-jets}, the extra tensor is a
smooth quadratic polynomial to leading order, with remainder
whose coefficients and derivatives through order two have size
$O(\sigma^{3-j})$.  Its second derivatives are Lipschitz away
from the axis and extend Hölder continuously with any exponent
less than one.  The tangential correction is
$(z_1^2+z_2^2)h_2+O(\sigma^3)$ and has the same property.
Thus $\mathfrak h$ is a positive $C^{2,\alpha}$ metric across
the disk center.  Its Ricci tensor vanishes off the center by
the static equations, and on the center by continuity.
Harmonic-coordinate elliptic regularity for the Ricci tensor
now makes this metric smooth: harmonic coordinates are
$C^{3,\alpha}$, the Ricci equation has principal part
$-\tfrac12\mathfrak h^{ab}\partial_{ab}\mathfrak h_{ij}$,
and repeated Schauder estimates apply to its quadratic
first-derivative terms \cite{DeTurckKazdan1981}.

The circle manifold is complete.  Project a finite-length curve
to the base, using Equation~\eqref{eq:static:circle}; the
preceding length test gives a base limit.  Over an interior
positive base limit the circle metric is uniformly
nondegenerate and the fiber is compact.  At $S$ the disk
coordinates give the limit.  These observations also prove
metric completeness by joining a summably Cauchy subsequence
with paths of summable lengths.

\smallskip\noindent
\textbf{Step 2: Excluding an interior frontier.}
Suppose an interior frontier of $\mathcal C$ remained.  A
sequence tending to it in the original metric escapes every
compact subset of the circle manifold: the continuous map
from the circle attachment to the original exterior maps its
compact subsets to compact sets avoiding this frontier.
Choose a connected product cross-section
$\{|x|=R\}\times S^1$ sufficiently far out.  Its outer side
is the asymptotically Euclidean base times a circle, and is
noncompact.  A connected two-sided cross-section has at most two
complement components: each component must meet one of its two
connected collar sides.  Thus the inner side is a single component
and contains the escaping sequence just constructed.  The
cross-section therefore separates at least two ends of a
complete smooth manifold with nonnegative Ricci curvature.
Minimizing segments between escaping sequences in the two
sides cross the compact separator and have a subsequence
converging to a globally minimizing line.  The
Cheeger--Gromoll splitting theorem
\cite{CheegerGromoll1971} gives an isometric product
$\R\times N$.  If $N$ were noncompact this product would
have one end: outside any compact cylinder its two axial
sides can be joined through a point of $N$ outside its compact
projection.  Hence $N$ is compact, and its ball volumes grow
at most linearly.  On the other hand, in the original product
end, $h$ tends after a linear change to the Euclidean metric
and $\lambda\to1$.  Uniform metric comparison places a
coordinate annulus of radius proportional to $L$, times the
entire circle, inside a ball of radius $L$ about a fixed
point; its volume is at least $c_0L^n$.  This contradicts
linear growth.  There is no interior frontier.  Since
$\Omega$ is connected, $\mathcal C=\operatorname{int}\Omega$.
\end{proof}

\subsection{The static end and its conformal compactification}
\label{subsec:static:compactification}

There are now no interior null frontiers, so the static base has a
compact core and exactly one end.  This is the point at which conformal
doubling becomes available.  The conformal joining and compactification
are the static uniqueness method of Bunting--Masood-ul-Alam
\cite{BuntingMasoodUlAlam1987}, developed in higher dimensions by
Gibbons--Ida--Shiromizu
\cite[Section~2, Equations~(2.6)--(2.12)]{GibbonsIdaShiromizu2002}.
We give the asymptotic calculation and the regularity of the added
point explicitly, rather than importing a uniqueness theorem with
different global hypotheses.

We now use the complete static base just constructed.  In this subsection
write $h=h_s$ and set $d=n-2=k-1$.  Thus $S$ is compact and connected,
$h$ is smooth up to $S$, and
\begin{equation}
 \Hess_h\lambda=\lambda\Ric_h,\qquad
 \Delta_h\lambda=0,\qquad R_h=0,\qquad
 \lambda|_S=0,\qquad \partial_{\nu_h}\lambda|_S=\kappa>0.
 \label{eq:static:base-for-double}
\end{equation}
Here $\nu_h$ is the $h$-unit normal pointing from $S$ into the static
exterior; on large coordinate spheres it denotes the $h$-unit normal
toward infinity.  We have $0<\lambda<1$ in the interior, $S$ is
totally geodesic, and the
complement of the one asymptotically Euclidean end is compact.  A fixed
linear change of end coordinates makes the limiting metric Euclidean.
We may decrease the decay exponent and choose
\[
 \frac d2<q_0<\min(q,d),\qquad q_0\ne1,
 \qquad h-\delta=O_2(r^{-q_0}),\qquad
 \lambda-1=O_2(r^{-q_0}).
\]

The required end expansion has two purposes: the retained end must
have zero mass, and the other end must become a metric with enough
Sobolev regularity at a single added point.  We first record the
potential estimate that isolates the leading monopole.

\begin{lemma}[The monopole of a decaying Poisson solution]
\label{lem:static:poisson-tail}
Suppose $w$ is smooth outside a ball in $\R^n$, tends to zero there,
and $\Delta_\delta w=O_j(r^{-n-\epsilon})$ for every $j\ge0$, where
$\epsilon>0$.  Suppose initially that $w=O(r^{-b})$ for some $b>0$.
For every $0<\eta<\min(1,\epsilon)$ there is a constant $A$ such that
\[
 w=A r^{-d}+O_j(r^{-d-\eta})\qquad(j\ge0).
\]
The assertion also holds componentwise for a tensor written in these
coordinates.
\end{lemma}

\begin{proof}
Extend $\chi w$ by zero across a ball, where $\chi$ is one far out, and
put $f=\Delta_\delta(\chi w)$.  Then
$\int(1+|z|^\eta)|f(z)|\,dz<\infty$.  For the fundamental solution
$\Phi(x)=-(d\omega)^{-1}|x|^{-d}$ of $\Delta_\delta$, the difference
$\chi w-\Phi*f$ is an entire harmonic function tending to zero, and
therefore vanishes.  In the region $|z|<|x|/2$, the estimate
\[
 |\Phi(x-z)-\Phi(x)|
 \le C|x|^{-d-\eta}|z|^\eta
\]
controls the difference from the monopole.  In the remaining region,
split off $|x-z|<|x|/2$.  On that ball, the pointwise bound for $f$
gives $C|x|^{-d-\epsilon}$ after integration against $\Phi$; on its
complement the finite $\eta$-moment gives $C|x|^{-d-\eta}$.
Consequently $A=-(d\omega)^{-1}\int f$ and the asserted zeroth-order
remainder follows.  Interior Poisson estimates on annuli rescaled to
unit size give every differentiated remainder.  This last step avoids
differentiating the singular Newton kernel twice under an absolutely
convergent integral.
\end{proof}

\begin{proposition}[Harmonic asymptotics of the static end]
\label{prop:static:harmonic-end}
There are end coordinates and a number
$0<\eta<\min(1,2q_0-d)$ such that, for every $j\ge0$,
\begin{equation}
 \begin{gathered}
 \check h:=\lambda^{2/d}h
       =\delta+O_j(r^{-d-\eta}),\qquad
 V:=\log\lambda=-a r^{-d}+O_j(r^{-d-\eta}),\\
 a=\frac{\kappa\Area_h(S)}{d\omega}>0.
 \end{gathered}
 \label{eq:static:harmonic-end}
\end{equation}
\end{proposition}

\begin{proof}
The conformal Ricci and Laplace formulas, with conformal exponent
$V/d$, give
\begin{equation}
 \Ric_{\check h}=\frac{k}{d}\,dV\otimes dV,
 \qquad \Delta_{\check h}V=0.
 \label{eq:static:optical-system}
\end{equation}
For example $\Delta_hV=-|dV|_h^2$ and
$\Ric_h=\Hess_hV+dV\otimes dV$; substitution cancels both the
Hessian and the scalar multiple of $h$ in the conformal Ricci formula.

Here are details of the choice of harmonic coordinates.  Extend the
end coefficients to be Euclidean inside a large radius $R$, using a
cutoff on $R<r<2R$, and solve for $y^i=x^i+\psi^i$ satisfying
$\Delta_{\check h}y^i=0$ on the original tail.  The weight for $\psi$
is $\gamma=1-q_0\in(2-n,1)\setminus\{0\}$.  The Euclidean Laplacian
has a bounded right inverse from weighted $C^{0,\alpha}_{\gamma-2}$
to weighted $C^{2,\alpha}_\gamma$.  To see the bound directly when
$\gamma<0$, use its Newton integral; when $\gamma>0$, replace the
kernel on $|z|>1$ by $\Phi(x-z)-\Phi(-z)$.  The latter subtraction
makes the outer integral convergent.  Splitting the integral into
$|z|<|x|/2$, comparable radii, and $|z|>2|x|$ gives growth
$(1+|x|)^\gamma$, and rescaled interior estimates give its derivatives
and H\"older seminorm.  The restrictions on $\gamma$ are exactly those
used in these integrals.  Constants in the positive-weight kernel
cause no difficulty; this construction fixes one right inverse.

In coordinates $x=Rz$, the extended coefficients differ from their
Euclidean values by $O(R^{-q_0})$ in the corresponding weighted
coefficient norms.  Hence
\[
 \psi=T\bigl(-\Delta_{\check h}x
                 -(\Delta_{\check h}-\Delta_\delta)\psi\bigr)
\]
is solved by a Neumann series for large $R$.  It gives
$\psi=O_2(r^{1-q_0})$, with the scaled H\"older control, and
$Dy=I+O(r^{-q_0})$.  Thus $y$ is a coordinate system sufficiently far
out: local invertibility follows from the last bound, and properness
and the estimate $y=x+o(r)$ give injectivity and coverage of a smaller
tail.  The initial $O_2$ coefficient bounds also give scaled
$W^{3,p}$ bounds for $y$ for every finite $p$ by differentiating its
scalar equation once.  The metric in $y$ consequently has controlled
$W^{2,p}$ and $C^{1,\alpha}$ norms on scaled annuli, for $p>n$.

In these coordinates the equations have the diagonal principal part
\[
 \check h^{ab}\partial_{ab}\check h_{ij}
   =Q_{ij}(\check h^{-1},D\check h)
      -\frac{2k}{d}\,\partial_iV\partial_jV,
 \qquad
 \check h^{ab}\partial_{ab}V=0,
\]
where $Q$ is quadratic in $D\check h$.  Scalar interior Schauder
estimates, first with the controlled $C^{1,\alpha}$ coefficients and
then after differentiation, give
$\check h-\delta,V=O_j(r^{-q_0})$ for every $j$.
It follows that
\[
 \Delta_\delta(\check h_{ij}-\delta_{ij}),\quad
 \Delta_\delta V=O_j(r^{-2-2q_0}).
\]
Since $2+2q_0>n$, Lemma~\ref{lem:static:poisson-tail} gives
\[
 \check h_{ij}=\delta_{ij}+A_{ij}r^{-d}
                  +O_j(r^{-d-\eta}),\qquad
 V=A_0r^{-d}+O_j(r^{-d-\eta}).
\]
The harmonic coordinate identity
$\partial_j(\sqrt{\det\check h}\,\check h^{ij})=0$ has leading term
\[
 d\left(A_{ij}-\tfrac12\tr(A)\delta_{ij}\right)x^j r^{-n}=0.
\]
All omitted terms are smaller by a positive power of $r$.  Letting
$r\to\infty$ on each ray shows
$A_{ij}-\tfrac12\tr(A)\delta_{ij}=0$.  Its trace is
$(1-n/2)\tr(A)=0$, so $A=0$ because $n\ge3$.

Write $A_0=-a$.  As $\eta<1\le d$, exponentiation gives
$\lambda=1-a r^{-d}+O_j(r^{-d-\eta})$.
Integrating $\Delta_h\lambda=0$ between $S$ and a large coordinate
sphere, with the inner integration normal $-\nu_h$, gives
\[
 \lim_{r\to\infty}\int_{S_r}\partial_{\nu_h}\lambda\,dA_h
   =\kappa\Area_h(S).
\]
The left side is $d\omega a$ by the expansion.  This proves the
coefficient and its strict positivity.
\end{proof}

We have determined the leading lapse coefficient and removed the
metric monopole in the conformal metric $\check h$.  These two facts
supply exactly the cancellation and point regularity needed for the
double.  Reflect $h$ across $S$ and extend $\lambda$ oddly.  The reflected
pair is $C^{2,\alpha}$ for every $0<\alpha<1$.  Indeed, in Gaussian
distance $\sigma\ge0$,
\[
 h=d\sigma^2+\gamma(\sigma),\qquad
 \gamma'(0)=0,\qquad
 \lambda=\kappa\sigma+O(\sigma^3).
\]
The first identity at $S$ is total geodesicity; the missing quadratic
lapse term follows from $\Hess_h\lambda=0$ there.  Even reflection of
$\gamma$ and odd reflection of $\lambda$ therefore have continuous
second derivatives and locally Lipschitz second derivatives in these
smooth one-sided charts.  The static equations and scalar flatness
hold across the seam by continuity.  On this double define
\begin{equation}
 \gamma_* =\left(\frac{1+\lambda}{2}\right)^{4/d}h.
 \label{eq:static:doubled-metric}
\end{equation}
On the second copy this is understood as
$((1-\lambda)/2)^{4/d}h$ with the original positive lapse.

\begin{lemma}[The filled double]
\label{lem:static:filled-double}
Adding one point at the second end produces a connected complete
boundaryless manifold $N$ with a locally uniformly elliptic metric
$\gamma_*$ of class $W^{2,s}_{\mathrm{loc}}$, for some $s>n/2$.
It is scalar flat in distributions, smooth except possibly at that
point and at the $C^{2,\alpha}$ seam, and has one smooth
asymptotically Euclidean end of ADM mass zero.
\end{lemma}

\begin{proof}
The scalar conformal formula and $\Delta_h\lambda=0$ make
Equation~\eqref{eq:static:doubled-metric} scalar flat away from the
added point, including the seam.  At the retained end its exact
expression relative to $\check h$ is
\[
 \gamma_*
 =\left(\frac{(1+\lambda)^2}{4\lambda}\right)^{2/d}\check h
 =\left(1+\frac{(1-\lambda)^2}{4\lambda}\right)^{2/d}\check h.
\]
Proposition~\ref{prop:static:harmonic-end} therefore gives
$\gamma_* =\delta+O_j(r^{-d-\eta})+O_j(r^{-2d})$.
Its ADM flux is zero.

On the second copy the factor relative to $\check h$ is
\[
 \left(\frac{1-\lambda}{2}\right)^{4/d}\lambda^{-2/d}
  =\left(\frac a2\right)^{4/d}r^{-4}
      \bigl(1+O_j(r^{-\eta})\bigr).
\]
In inverted coordinates $z=x/|x|^2$, $\rho=|z|$, the identity
$r^{-4}\delta_{ij}\,dx^i dx^j=\delta_{ij}\,dz^i dz^j$ gives
\begin{equation}
 (\gamma_*)_{ij}(z)
   =\left(\frac a2\right)^{4/d}
          \bigl(\delta_{ij}+e_{ij}(z)\bigr),
 \qquad |D^j e(z)|\le C_j\rho^{\eta-j}\quad(j\ge0).
 \label{eq:static:inverted-point}
\end{equation}
The orthogonal reflection in the differential of inversion causes
no change in these orders.  Define $e(0)=0$.  Choose
\begin{equation}
 \frac n2<s<\frac{n}{2-\eta}<n.
 \label{eq:static:point-exponents}
\end{equation}
Then $\int_0^1\rho^{n-1+s(\eta-2)}\,d\rho<\infty$.
Thus the metric has two weak derivatives in $L^s$ across the point.
For completeness, integrating by parts over $\rho>\varepsilon$
introduces no missing second-derivative measure: the potentially
relevant first-derivative boundary term is
$O(\varepsilon^{n+\eta-2})\to0$; the zeroth-order boundary term also
vanishes.  Its curvature is consequently the ordinary $L^s$ tensor
computed from these weak derivatives.  Indeed, with
\[
 p=\frac{ns}{n-s}>n,\qquad p>2s,
\]
the first derivatives belong to $L^p$, so their quadratic products
belong to $L^{p/2}\subset L^s$.  Scalar curvature vanishes almost
everywhere and therefore distributionally everywhere.

Equation~\eqref{eq:static:inverted-point} gives a continuous positive
definite extension at the added point.  The remainder of the compact
core has the same property.  Every escape from this core along the
retained end has infinite length, since its metric is uniformly
comparable to the Euclidean metric.  This proves completeness of the
filled metric.
\end{proof}

\subsection{Rigidity of the zero-mass double}
\label{subsec:static:zero-mass}

The metric $\gamma_*$ is complete and scalar flat, has a single
asymptotically Euclidean end of mass zero, and is $W^{2,s}_{\mathrm{loc}}$
for $n/2<s<n$.  Its possible nonsmoothness is confined to the joining
hypersurface and the filled point.  We will show that it is Euclidean.
The main step is to vary this metric without leaving the scalar-flat
class: nonnegative mass for both signs of a compact variation will
force its Ricci tensor to vanish.

The numerical Riemannian Penrose inequality in
Theorem~\ref{thm:riemannian-input} is the input.  It is supplied by
\emph{Conformal flow and the Riemannian Penrose inequality with
minimizing frontiers}.  We first construct the Laplace inverse needed
both to insert a Green-function pole and to correct scalar curvature.
Enclosing the pole by a minimizing frontier gives smooth nonnegative
mass.  We then smooth compact Sobolev defects and correct the scalar
curvature, while computing the change in mass.  Finally we remove the
weak regularity defects of the Ricci-flat double and use volume
comparison.  Only the numerical
inequality is needed from the Riemannian theorem.

Throughout this subsection $n\ge3$, $d=n-2$, $k=n-1$, and
$\omega=|S^{n-1}|$.  The Laplacian is $\Delta=\operatorname{div}\nabla$,
and the ADM mass has factor $1/(2k\omega)$.

\begin{theorem}[Riemannian numerical input {\cite[Theorem~1.1]{CompanionRiemannian}}]
\label{thm:riemannian-input}
Let $n\ge3$, $d=n-2$, $k=n-1$, and $\omega=|S^{n-1}|$.
Let $(E,h)$ be a connected complete \emph{enclosing-set exterior}:
the end-containing closed region outside a filled set, with its
compact full frontier $\Sigma$ included.  For $p,q\in E$, define the
possibly extended rectifiable-path distance
\[
 d_E(p,q)=\inf\{\operatorname{Length}_h(\gamma):
       \gamma\subset E\text{ is a rectifiable path joining }p\text{ to }q\},
 \qquad \inf\varnothing=+\infty.
\]
Completeness means that every finite-distance equivalence class
$\{q\in E:d_E(p,q)<\infty\}$ is complete for the restricted metric.
All points and components of $\Sigma$ remain in $E$, and their full
perimeter is counted.  This convention neither identifies $E$ with
the length completion of its open exterior nor asserts finite-length
access from the end to every singular frontier point.
Assume that the ambient metric $h$ is smooth on $E$ and through
$\Sigma$, and that the following conditions hold:
\begin{enumerate}
\item There is one Euclidean coordinate end with compact complement,
      and $h-\delta=O_2(r^{-\tau})$ for some $\tau>d/2$.
\item The scalar curvature satisfies $R_h\ge0$,
      $R_h\in L^1(E,dV_h)$, and $R_h=O(r^{-n-\beta})$ on the end
      for some $\beta>0$.
\item The full frontier $\Sigma$ is outer minimizing and locally
      perimeter minimizing as the boundary of the filled side.
      Its regular part is a smooth embedded minimal hypersurface,
      and its singular set has Hausdorff dimension at most $n-8$.
\end{enumerate}
Then, with $A$ the reduced-boundary perimeter of $\Sigma$, counting
every component,
\[
 m_{\rm ADM}(h)\ge \frac12\left(\frac{A}{\omega}\right)^{d/k}.
\]
The assertion includes $A=0$ and requires neither a spin assumption
nor connectedness of $\Sigma$.
\end{theorem}

Below we construct a nonempty detached
frontier in a smooth scalar-flat metric and verify these hypotheses
before applying the inequality.

\begin{lemma}[A compact-source Laplace inverse]
\label{lem:static:compact-source}
Let $\gamma$ be a complete locally uniformly elliptic
$W^{2,s}_{\mathrm{loc}}$ metric on a connected boundaryless manifold,
where $n/2<s<n$.  Suppose that it has compact complement to one
smooth end on which $\gamma-\delta=O_j(r^{-d})$ for every $j$.
Fix a compact set $K$ and $d/2<b<d$.  For every $f\in L^s$ supported
in $K$ there is a unique weak solution $v\in W^{2,s}_{\mathrm{loc}}$
of $\Delta_\gamma v=f$ tending to zero at the end.  It satisfies
\begin{equation}
 \|v\|_{W^{2,s}(K_1)}+
 \sup_{r\ge R}r^b\sum_{j=0}^2r^j
          \|D^jv\|_{L^\infty(\{r<|x|<2r\})}
       \le C\|f\|_{L^s(K)},
 \label{eq:static:inverse-estimate}
\end{equation}
where $K_1$ contains $K$ and the core.  Moreover
\begin{equation}
 v=A_f r^{-d}+O_j(r^{-d-\eta_f}),\qquad
 A_f=-\frac{1}{d\omega}\int_N f\,dV_\gamma,
 \label{eq:static:source-flux}
\end{equation}
for every fixed $0<\eta_f<\min\{1,b\}$.  The inverse bound is uniform
under sufficiently small $W^{2,s}$ perturbations supported in a fixed
compact set; the end is kept unchanged.
\end{lemma}

\begin{proof}
We describe the underlying global estimate to fix the solvability
class.  This manifold has the Euclidean Sobolev inequality
\begin{equation}
 \|\varphi\|_{L^{2n/(n-2)}}
       \le C\|d\varphi\|_{L^2},\qquad
       \varphi\in C_c^\infty(N).
 \label{eq:static:global-sobolev}
\end{equation}
To verify the compact part of this assertion, on the end write
\[
 \varphi(r,\theta)=-\int_r^\infty\partial_t\varphi(t,\theta)\,dt.
\]
Weighted Cauchy--Schwarz gives
\[
 |\varphi(r,\theta)|^2
 \le\frac{r^{-d}}d
       \int_r^\infty|\partial_t\varphi|^2t^{n-1}\,dt.
\]
Integration over a fixed end annulus bounds its $L^2$ norm by the
global Dirichlet integral.  A Poincar\'e inequality on a connected
compact region meeting this annulus gives the same bound on the
core.  Finite-chart Sobolev inequalities on the core, and the
Euclidean Sobolev inequality applied after an end cutoff, now prove
Equation~\eqref{eq:static:global-sobolev}.  Uniform metric
comparability suffices in this argument.

Let $\mathcal D^{1,2}$ be the completion of $C_c^\infty(N)$ in the
Dirichlet norm.  Since a compactly supported $L^s$ function belongs
to $L^{2n/(n+2)}$, Equation~\eqref{eq:static:global-sobolev} and the
Lax--Milgram theorem solve
\[
 \int_N\langle dv,d\varphi\rangle_\gamma\,dV_\gamma
       =-\int_N f\varphi\,dV_\gamma,
 \qquad v\in\mathcal D^{1,2}.
\]
They also bound $\|dv\|_2$ and $\|v\|_{2n/(n-2)}$ by $C\|f\|_s$.
Local Moser iteration for this divergence equation gives a local
$L^\infty$ bound, because $s>n/2$.  The local $W^{2,s}$ estimate
applies as follows.  In coordinates write
$\Delta_\gamma=a^{ij}\partial_{ij}+B^j\partial_j$.
Here $a$ is continuous and uniformly elliptic, and
$B\in L^p$ with $p=ns/(n-s)>n$.  On small balls, freeze $a$, use the
constant-coefficient Calder\'on--Zygmund estimate, and absorb its
oscillation.  The drift is absorbed by H\"older and interpolation,
\[
 \|B Dv\|_s\le\|B\|_p\|Dv\|_n,
 \qquad
 \|Dv\|_n\le\varepsilon\|D^2v\|_s+C_\varepsilon\|v\|_s,
\]
where the strict inequality $p>n$ permits interpolation below the
Sobolev exponent for $Dv$.  These scalar interior estimates can
first be applied to smooth metric and source approximations; their
uniform energy, Moser, and interior bounds pass to the weak solution.
This proves the core part of
Equation~\eqref{eq:static:inverse-estimate}; the scalar estimates used
here are those of~\cite{GilbargTrudinger2001}.  The compact bound controls
the inner boundary of the end.  We next turn it into the decay and
flux estimates that will determine the conformal mass change.

For sufficiently large $r$,
\[
 \Delta_\gamma r^{-b}
      =b(b-d)r^{-b-2}+O(r^{-b-d-2})<0.
\]
Compare $v$ and $-v$ with $Cr^{-b}$, choosing $C$ from the fixed
inner sphere bound.  The comparison is legitimate on the unbounded
end by testing with $(\pm v-Cr^{-b})_+$: this function has zero inner
trace and belongs to $\mathcal D^{1,2}$ since $2b>d$.
It follows that $|v|\le Cr^{-b}$.  Rescaled interior estimates for
the smooth homogeneous end equation give all derivative bounds
there.  A decaying homogeneous solution has neither a positive
maximum nor a negative minimum by the weak maximum principle, so
the inverse is unique in the stated class.

Finally,
$\Delta_\delta v=O_j(r^{-n-b})$ on the end.  Apply
Lemma~\ref{lem:static:poisson-tail} and integrate the weak equation
over large coordinate balls to obtain
\[
 -d\omega A_f
   =\lim_{r\to\infty}\int_{S_r}\partial_{\nu_\gamma}v\,dA_\gamma
   =\int_N f\,dV_\gamma.
\]
There is no interior boundary in this distributional identity.
All constants above stay bounded under small compact $W^{2,s}$
metric perturbations: their $C^0$ norms and ellipticity, H\"older
moduli of $a$, and $L^p$ norms of $B$ are uniformly controlled, and
the end is fixed.
\end{proof}

\begin{lemma}[Smooth nonnegativity from the numerical Penrose theorem]
\label{lem:static:smooth-nonnegative-mass}
Assume the numerical assertion of Theorem~\ref{thm:riemannian-input}.
Every complete smooth connected boundaryless scalar-flat manifold
with compact complement to one end satisfying
$\gamma-\delta=O_j(r^{-d})$ has nonnegative ADM mass.
\end{lemma}

\begin{proof}
Fix a smooth point $p$.  A local Dirichlet Green function, cut off
inside its coordinate ball, has distributional Laplacian
$-d\omega\delta_p+f_0$, where $f_0$ is smooth and supported away
from $p$.  Subtract the solution of
$\Delta_\gamma v=f_0$ supplied by
Lemma~\ref{lem:static:compact-source}.  The resulting function $G$
is harmonic off $p$, tends to zero at the end, and has the positive
pole $G=r_p^{-d}(1+o(1))$ in geodesic distance from $p$, with the
corresponding differentiated leading term.  The maximum principle
on the punctured manifold gives $G>0$.  Its distributional equation
and the same flux calculation as in
Equation~\eqref{eq:static:source-flux} give
\[
 \Delta_\gamma G=-d\omega\delta_p,
 \qquad G=r^{-d}+O_j(r^{-d-\eta_G})\quad\hbox{at infinity}.
\]
For $\varepsilon>0$ put
$\gamma_\varepsilon=(1+\varepsilon G)^{4/d}\gamma$ on
$N\setminus\{p\}$.  This is scalar flat and complete, with the
original end and a second complete end at the pole.  Near the pole,
the length factor is comparable to $r_p^{-2}$.

For fixed $\varepsilon$, sufficiently small spheres about $p$ have
strictly negative mean curvature toward the original end.  In fact,
writing $U=1+\varepsilon G$, the conformal mean-curvature law is
\[
 H_{\gamma_\varepsilon}
 =U^{-2/d}\left(H_\gamma+\frac{2k}{d}\partial_{\nu_\gamma}\log U\right),
\]
and the expression in parentheses is
$k/r_p-2k/r_p+o(r_p^{-1})<0$.
Choose one such sphere $S_\rho$.  Truncate the pole end at a
smaller sphere $S_{\rho_0}$, where $0<\rho_0<\rho$, retaining the
closed ambient exterior $M_{\rho_0}$.  This is a smooth scalar-flat
manifold with compact boundary and compact complement to the original
Euclidean end.  Its ordinary metric completeness follows from its
compact smooth core, smooth boundary charts, and complete
asymptotically Euclidean tail.  This ambient completeness is independent
of paths in the singular enclosure to be constructed.  We keep
$M_{\rho_0}$ as an ambient neighborhood on both sides of the minimizing
frontier that will be constructed.

For the enclosure construction itself, let $X_\rho$ be the original-end
exterior bounded by $B=S_\rho$, with $B$ included.  Apply the smooth
full-enclosure results of
\cite[Lemmas~2.1--2.3]{CompanionRiemannian} to
$(X_\rho,\gamma_\varepsilon)$.  These results assume a smooth connected
manifold complete with its nonempty compact smooth boundary included,
one Euclidean coordinate end with compact complement, and
$g-\delta=O_2(r^{-q})$ for some $q>0$.  A compact smooth filling $O$
is attached behind $B$, and perimeter is measured in the resulting
boundaryless manifold, counting all free-frontier components and any
contact with $B$.  In our application choose the metric on that
auxiliary filling to agree with $\gamma_\varepsilon$ in a collar
behind $S_\rho$, where the original metric is smooth.  Its extension
farther inside the filling has no role in the scalar-curvature
argument on the exterior.

The end condition holds because
$\gamma_\varepsilon-\delta=O_j(r^{-d})$; in particular, large
coordinate spheres have mean curvature $k/r+O(r^{-d-1})>0$.
The enclosure lemmas give a compact full-perimeter minimizer with
connected end exterior.  Since $H_B<0$ toward $X_\rho$, their
strict-barrier conclusion puts a fixed exterior collar of $B$ inside
the minimizing filled set.  Its frontier is therefore detached from
$B$, is outer minimizing, and is locally perimeter minimizing in the
smooth metric.  Its singular set has Hausdorff dimension at most
$n-8$.  These statements use only smooth obstacle geometry and full
perimeter, independently of original-data variation or equality.

Consequently this compact frontier is locally perimeter minimizing,
outer minimizing, and minimal, with singular set of Hausdorff
dimension at most $n-8$ by minimizing-boundary regularity
\cite{Simon2018}.  Its area is positive, and the smooth metric is
defined on a neighborhood of it.  Its original-end exterior has
nonnegative scalar curvature, integrable scalar curvature (in fact
zero), and the decay $d>d/2$ required by
Theorem~\ref{thm:riemannian-input}.  To give its enclosing-set ambient realization,
use the already constructed $M_{\rho_0}$: its compact boundary is
strictly behind the detached frontier and its metric is smooth through
that frontier.  Let $D$ denote the retained closed enclosing-set
exterior, including its entire frontier, and let $d_D$ be the possibly
extended path distance defined above, using $\gamma_\varepsilon$.
We verify completeness of every finite-distance class of $d_D$.

Fix such a class $C$ and a $d_D$-Cauchy sequence $(x_j)$ in $C$.
Choose a subsequence $(x_{j_\ell})$ with
$d_D(x_{j_\ell},x_{j_{\ell+1}})<2^{-\ell-1}$.  By the definition
of the infimum, these successive points can be joined by rectifiable
paths $c_\ell\subset D$ of length less than $2^{-\ell}$.
Concatenate the paths.  Their total length is finite and their tail
lengths tend to zero.  Ambient distances between points on a tail are
bounded by that tail's length.  Ordinary completeness of $M_{\rho_0}$
therefore gives an ambient endpoint $x_\infty$, and $x_\infty\in D$
because $D$ is closed.  Appending this endpoint gives a rectifiable
path in $D$ with length at most the sum of the constituent lengths.
The appended limit lies in the same finite-distance class: the
completed concatenation joins it to the fixed chain point $x_{j_1}$
with finite length.  Its tails give
\[
 d_D(x_{j_\ell},x_\infty)
 \le\sum_{r\ge\ell}\operatorname{Length}_{\gamma_\varepsilon}(c_r)
 \longrightarrow0.
\]
The Cauchy property and the triangle inequality then give
$d_D(x_j,x_\infty)\to0$ for the whole sequence.  Thus $C$ is complete.
No restriction of $D$ to the finite-distance class of its end has been
made; all frontier points and components, and their full perimeter,
remain included.  The argument asserts neither finite-length access
to an arbitrary singular frontier point nor an identification with
the length completion of the open exterior.  Thus the numerical assertion of
Theorem~\ref{thm:riemannian-input} implies
$m_{\rm ADM}(\gamma_\varepsilon)\ge0$.
The old-end expansion of $G$ and the ADM normalization give
\[
 m_{\rm ADM}(\gamma_\varepsilon)
      =m_{\rm ADM}(\gamma)+2\varepsilon.
\]
Letting $\varepsilon\downarrow0$ proves the claim.  No property of
the discarded pole end other than the displayed barriers is used.
\end{proof}

\begin{lemma}[Conformal correction and mass at compact Sobolev defects]
\label{lem:static:rough-nonnegative-mass}
Under the numerical assertion of Theorem~\ref{thm:riemannian-input}, let $\gamma$
satisfy the hypotheses of Lemma~\ref{lem:static:compact-source} and
be scalar flat in distributions.  Suppose that it is smooth outside
a compact set.  Then $m_{\rm ADM}(\gamma)\ge0$.
For a sufficiently small compact $W^{2,s}$ metric perturbation
$\gamma'$ there is a positive factor $U=1+v$, tending to one,
such that $U^{4/d}\gamma'$ is scalar flat.  Its mass is
\begin{equation}
 m_{\rm ADM}(U^{4/d}\gamma')
   =m_{\rm ADM}(\gamma')
     -\frac{1}{2k\omega}\int_N R_{\gamma'}U\,dV_{\gamma'}.
 \label{eq:static:conformal-mass}
\end{equation}
The factor, its end monopole, and this mass depend differentiably
on a differentiable $W^{2,s}$ curve of perturbations supported in a
fixed compact set.
\end{lemma}

\begin{proof}
Put $c_n=4k/d$.  The conformal equation is
\[
 -c_n\Delta_{\gamma'}(1+v)+R_{\gamma'}(1+v)=0.
\]
Use the inverse $T_\gamma$ of
Lemma~\ref{lem:static:compact-source} to write it as
\begin{equation}
 v=T_\gamma\left[
       c_n^{-1}R_{\gamma'}(1+v)
          -(\Delta_{\gamma'}-\Delta_\gamma)v\right].
 \label{eq:static:conformal-neumann}
\end{equation}
The perturbation is supported in a fixed compact set $K$ containing
the nonsmooth locus.  Because $R_\gamma=0$, both terms in brackets
are supported in $K$.  The source norm is $L^s$ and the solution norm is that in
Equation~\eqref{eq:static:inverse-estimate}.  These products are
bounded in the indicated spaces: a second-order coefficient
perturbation is small in $L^\infty$, drift perturbations are small
in $L^p$, $Dv\in L^p$, $p/2>s$, and $v\in L^\infty$.
Also $R_{\gamma'}$ is small in $L^s$ because
$\gamma'\to\gamma$ in $W^{2,s}$ and $R_\gamma=0$.
Thus Equation~\eqref{eq:static:conformal-neumann} is solved by a
Neumann series, with $\|v\|\to0$ as $\gamma'\to\gamma$.
In particular $U>0$, the corrected metric is complete, and its
scalar curvature vanishes distributionally.  The coefficient maps $\gamma\mapsto\gamma^{-1}$,
$\gamma\mapsto\Delta_\gamma$, and $\gamma\mapsto R_\gamma$
are differentiable between these spaces: their derivatives contain
second derivatives linearly and products of two $L^p$ first
derivatives, with $p/2>s$.  Differentiating the convergent operator
inverse in Equation~\eqref{eq:static:conformal-neumann} therefore
proves differentiability in the perturbation parameter.

On the common end $v$ is harmonic for $\gamma'$, so it has the
monopole expansion of
Equation~\eqref{eq:static:source-flux}.  The ADM calculation for
$U=1+A r^{-d}+o_1(r^{-d})$ gives a mass change $2A$, or equivalently
\[
 m_{\rm ADM}(U^{4/d}\gamma')-m_{\rm ADM}(\gamma')
  =-\frac{2}{d\omega}
       \lim_{r\to\infty}\int_{S_r}
                   \partial_{\nu_{\gamma'}}U\,dA_{\gamma'}.
\]
Integrating $\Delta_{\gamma'}U=c_n^{-1}R_{\gamma'}U$ proves
Equation~\eqref{eq:static:conformal-mass}.  The flux integral is
finite, and its compact-source expression also proves its
differentiability.

To prove nonnegativity for the original rough metric, smooth its
coefficients in finitely many charts covering its nonsmooth compact
set, with a partition of unity, leaving the end fixed.  The resulting
smooth positive definite metrics $\gamma_j$ converge in $W^{2,s}$
on the core and uniformly; their scalar curvatures tend to zero in
$L^s$.  The preceding construction provides $U_j\to1$ and smooth
complete scalar-flat metrics $U_j^{4/d}\gamma_j$.  Smoothness of
$U_j$ follows from its smooth scalar elliptic equation.  The end
expansion is $\delta+O_\ell(r^{-d})$ at every derivative order $\ell$, so
Lemma~\ref{lem:static:smooth-nonnegative-mass} applies.  Compact
smoothing leaves the uncorrected end mass unchanged, and
Equation~\eqref{eq:static:conformal-mass} gives
\[
 0\le m_{\rm ADM}(U_j^{4/d}\gamma_j)
          \longrightarrow m_{\rm ADM}(\gamma).
\]
This proves nonnegativity at the asserted regularity as well.
\end{proof}

We now have nonnegative mass at the precise regularity of the filled
double and under its compact conformally corrected perturbations.
Applying this statement to both signs of a variation converts its
zero mass into Ricci flatness before any smooth rigidity theorem is
used.

\begin{proposition}[Euclidean conformal double]
\label{prop:static:euclidean-double}
The filled metric $(N,\gamma_*)$ of
Lemma~\ref{lem:static:filled-double} is isometric to Euclidean
$\R^n$.
\end{proposition}

\begin{proof}
\smallskip\noindent
\textbf{Step 1: Mass variation and weak Ricci flatness.}
Let $v_{ij}$ be an arbitrary smooth symmetric tensor compactly
supported in a smooth patch of $\gamma_*$, and put
$\gamma_t=\gamma_*+tv$.  For both signs of sufficiently small $t$,
Lemma~\ref{lem:static:rough-nonnegative-mass} constructs a positive
factor $U_t$ such that $\widetilde\gamma_t=U_t^{4/d}\gamma_t$ is
scalar flat.  The same lemma gives
$m_{\rm ADM}(\widetilde\gamma_t)\ge0$.
Since $U_0=1$ and $m_{\rm ADM}(\gamma_*)=0$, this differentiable
mass has derivative zero at $t=0$.
Equation~\eqref{eq:static:conformal-mass} gives precisely
\begin{align*}
 0
 &= -\frac{1}{2k\omega}\int_N
          R'_{\gamma_*}[v]\,dV_{\gamma_*}\\
 &= \frac{1}{2k\omega}\int_N
          \langle\Ric_{\gamma_*},v\rangle_{\gamma_*}\,dV_{\gamma_*}.
\end{align*}
For the second equality use
$R'[v]=\divg\divg v-\Delta\tr v-\langle\Ric,v\rangle$;
the divergence terms integrate to zero because the test is
supported in that smooth patch.  Hence $\Ric_{\gamma_*}=0$ on every
smooth patch.  Its coordinate components belong to $L^s$, as proved
in Lemma~\ref{lem:static:filled-double}.  The seam and the filled
point have zero volume, so $\Ric_{\gamma_*}=0$ also as a
distribution everywhere.

\smallskip\noindent
\textbf{Step 2: Smoothness at the seam and the filled point.}
We include the regularity argument before applying a smooth
comparison theorem.  For a locally uniformly elliptic
$W^{2,s}$ metric with $n/2<s<n$, put $p=ns/(n-s)>n$.
It belongs to $W^{1,p}$, so local harmonic coordinates exist as
$C^1$ diffeomorphisms of class $W^{2,p}$ by
\cite[Corollary~2.8]{JulinLiimatainenSalo2017} with harmonic exponent
two.  In original coordinates their equation is
\[
 a^{ij}\partial_{ij}y^\ell+B^j\partial_jy^\ell=0,
 \qquad a\in W^{2,s},\quad B\in W^{1,s}\cap L^p.
\]
We justify the extra derivative without assuming strong regularity
of the differentiated coordinates.  Define the distributional operator
on $W^{1,p}_{\mathrm{loc}}$ functions by
\[
 Lz=\partial_i(a^{ij}\partial_jz)
           +(B^j-\partial_i a^{ij})\partial_jz .
\]
Its lower-order product belongs to $L^{p/2}\subset L^s$ locally;
on $W^{2,s}$ functions it equals
$a^{ij}\partial_{ij}z+B^j\partial_jz$.
For $z=\partial_k y^\ell\in W^{1,p}$, differentiating the scalar
coordinate equation by weak product rules therefore gives
\[
 Lz=F,\qquad
 F=-(\partial_k a^{ij})\partial_{ij}y^\ell
            -(\partial_kB^j)\partial_jy^\ell\in L^s.
\]
Indeed $Dy$ is bounded, $DB\in L^s$, and
$Da,D^2y\in L^p$ with $p/2>s$.

Choose a small coordinate ball $B_r$ and a cutoff
$\chi\in C_c^\infty(B_r)$ equal to one on a smaller ball.
Then $w=\chi z\in W^{1,p}_0(B_r)$ satisfies $Lw=H\in L^s(B_r)$:
the additional terms are
$2a^{ij}(\partial_i\chi)(\partial_jz)$ and
$(a^{ij}\partial_{ij}\chi+B^j\partial_j\chi)z$.
Construct $v\in W^{2,s}(B_r)\cap W^{1,s}_0(B_r)$ with $Lv=H$
by perturbing the constant-coefficient Dirichlet inverse with
coefficient matrix $a(x_0)$ at the ball center.  Its
second-derivative estimate and the scaled Sobolev inequality give
\[
 \|B Dv\|_{L^s(B_r)}
 \le C r^{\,1-n/p}\|B\|_{L^p(B_r)}
                  \|D^2v\|_{L^s(B_r)},\qquad
 1-\frac np=2-\frac ns>0 .
\]
Continuity of $a$ makes the other operator error at most
$C\|a-a(x_0)\|_\infty\|D^2v\|_s$.  Both errors are small
for sufficiently small $r$, so a Neumann series constructs $v$.

Sobolev embedding gives $v\in W^{1,p}$ with zero trace.  Thus
$q=w-v\in W^{1,p}_0(B_r)\subset H^1_0(B_r)$ satisfies
\[
 \partial_i(a^{ij}\partial_jq)
          +(B^j-\partial_i a^{ij})\partial_jq=0 .
\]
Testing with $q$, using ellipticity with constant $\lambda_0>0$
and the Euclidean Sobolev inequality, gives
\[
 \lambda_0\|Dq\|_2^2
 \le C\|B-\operatorname{div}a\|_{L^n(B_r)}\|Dq\|_2^2 .
\]
The coefficient norm is arbitrarily small on a small ball because
$B,Da\in L^p$ and $p>n$.  Hence $q=0$.  On the smaller ball
$z=v\in W^{2,s}$, proving $y\in W^{3,s}_{\mathrm{loc}}$.

The inverse coordinates have the same regularity: their third
derivatives contain $D^3y\in L^s$ and quadratic $D^2y$ terms in
$L^{p/2}\subset L^s$.  The metric transformation formula likewise
has highest terms $D^2\gamma_*$, $D^3(y^{-1})$, and products of
two $L^p$ factors.  Consequently the harmonic-coordinate metric
remains $W^{2,s}$.  To justify tensoriality, approximate the metric
in $W^{2,s}$ and the coordinate map in $W^{3,s}$.  The latter
convergence is also $C^1$, so the maps and their inverses remain
diffeomorphisms on smaller common charts.  Their transformed
metrics converge in $W^{2,s}$; the same product bounds imply
convergence of their Ricci tensors in $L^s$.  The smooth Ricci
transformation law therefore passes to the limit.  Composition
of the $L^s$ tensors causes no additional difficulty: approximate
the limiting tensor by a smooth tensor and use the uniformly
bounded Jacobians of the coordinate maps.

The harmonic-coordinate Ricci identity
\cite[Equation~(4.3)]{DeTurckKazdan1981} now gives the weak system
\[
 (\gamma_*)^{ab}\partial_{ab}(\gamma_*)_{ij}
       =Q_{ij}(\gamma_*^{-1},D\gamma_*),
\]
whose right side belongs to
$L^{s_1}$, where $s_1=ns/(2(n-s))>s$.
The principal coefficients are continuous and uniformly elliptic.
To obtain the higher regularity without assuming it in advance,
localize a component with a cutoff in a small harmonic-coordinate
ball.  Its equation has right side in $L^{s_1}$: besides the quadratic
term, the cutoff terms contain only the metric and its first
derivatives.  Solve that Dirichlet equation in $W^{2,s_1}$ by the
constant-coefficient inverse and the small oscillation of the
principal coefficients.  The difference from the original localized
component lies in $H^1_0$ and solves the homogeneous equation.
Writing the equation in divergence form introduces the drift
$-\partial_i(\gamma_*)^{ij}\in L^p$; its $L^n$ norm is small on
this ball.  The energy uniqueness argument just used for the
coordinate derivatives forces the difference to vanish.  Hence each
metric component belongs to $W^{2,s_1}$ on a smaller ball.
Repeating strictly improves the exponent until it exceeds $n$ in
finitely many steps.  Indeed while $s_j<n$ the reciprocal recurrence
is $n/s_{j+1}=2n/s_j-2$, starting below two.  If an exponent equals
$n$, the embedding into $W^{1,p'}$ for every finite $p'$ supplies
the next improvement above $n$.  Thus the metric becomes
$C^{1,\alpha}$; the displayed system and successive Schauder
estimates make it smooth.  The harmonic transition maps are then
smooth by their scalar harmonic equations.  This establishes a
smooth complete Ricci-flat Riemannian structure on the same metric
space, including the filled point and the seam.

\smallskip\noindent
\textbf{Step 3: Euclidean volume rigidity.}
The one asymptotically Euclidean end gives asymptotic volume ratio
one.  More explicitly, outside each sufficiently large fixed core
the metric lies between $(1-\epsilon)\delta$ and
$(1+\epsilon)\delta$, while the core contributes only a fixed
additive distance and finite volume.  Squeezing geodesic balls
between the corresponding coordinate balls and then sending
$\epsilon\downarrow0$ gives
\[
 \lim_{R\to\infty}
 \frac{\Vol_{\gamma_*}(B_{\gamma_*}(o,R))}{(\omega/n)R^n}=1.
\]
For a complete smooth manifold with nonnegative Ricci curvature,
Bishop--Gromov comparison makes this ratio nonincreasing; its
small-radius limit is also one.  The classical equality case of Bishop--Gromov comparison gives a
global Euclidean isometry; this volume-rigidity statement is recalled
in~\cite[p.~3237, before Theorem~1.1]{CavallettiManini}.
\end{proof}




\subsection{Recovering the round horizon and the original slice}

The double has now been identified with Euclidean space.  We finish
by locating its seam, solving an exterior harmonic Dirichlet problem,
and recovering the time coordinate of the original slice.  The final
boundary check is made in regular spacetime coordinates, because the
static time itself diverges on a future-horizon section.

We now finish the proof of Theorem~\ref{thm:static:classification}.
The preceding conformal-double argument identifies the complete
compactified double with Euclidean space.  On the retained side
its metric is
\begin{equation}\label{eq:static:euclidean-side}
 \delta=\left(\frac{1+\lambda}{2}\right)^{4/(n-2)}h.
\end{equation}
The seam is a compact connected embedded hypersurface.  Because
it was totally geodesic for $h$, the conformal second-form formula
makes it totally umbilic for $\delta$.  Its principal curvature
with the Euclidean unit normal $\nu_\delta$ toward the retained side is the positive constant
$2^{2/(n-2)}\,2\kappa/(n-2)$.
For completeness, a Euclidean umbilic hypersurface of dimension
$k\geq2$ is a sphere or a plane: Codazzi applied to $\mathrm{II}=(H/k)\delta|_{TS}$
forces $dH=0$; if its common principal curvature is $a\ne0$,
the ambient derivative of $p-a^{-1}\nu_\delta$ along every tangent
vector is zero.  It lies on the sphere with that center.  The
inclusion into that sphere is a local diffeomorphism with closed
and open image, and embeddedness gives the entire sphere.
The zero-curvature case would be an open subset of an affine
plane and cannot be compact.  Thus, after translation, the seam
is $\{|y|=R_0\}$, and the retained side is $\{|y|\geq R_0\}$
because it contains the unbounded end.

Let $\rho=|y|$ and $\Psi=2/(1+\lambda)$.  Equation~\eqref{eq:static:euclidean-side}
and $R_h=0$ give $\Delta_\delta\Psi=0$.  The boundary values
are $\Psi=2$ on $\rho=R_0$ and $\Psi\to1$ at infinity.
Comparison on bounded annuli followed by the limit at infinity
proves uniqueness of this exterior Dirichlet problem.  Hence
\begin{equation}\label{eq:static:isotropic-pair}
 \Psi=1+(R_0/\rho)^{n-2},\qquad
 h=\bigl(1+(R_0/\rho)^{n-2}\bigr)^{4/(n-2)}\delta,
 \qquad
 \lambda=\frac{1-(R_0/\rho)^{n-2}}
                  {1+(R_0/\rho)^{n-2}}.
\end{equation}
The areal radius is
$r=\rho(1+(R_0/\rho)^{n-2})^{2/(n-2)}$.
Substitution into Equation~\eqref{eq:static:isotropic-pair} gives
\[
 h=(1-2M/r^{n-2})^{-1}dr^2+r^2\sigma_k,
 \qquad \lambda^2=1-2M/r^{n-2},\qquad M=2R_0^{n-2}.
\]
At the seam $r=2^{2/(n-2)}R_0$ and its $h$-area equals its
original $g$-area.  Thus
$A=\omega(2M)^{k/(k-1)}$, and
Equation~\eqref{eq:static:constants} forces $M=m$.
This proves the identification of the static pair without
assuming topology or a vacuum development beforehand.

The identified base is the exterior of a ball in $\R^n$.
For $n\geq3$ it is simply connected.  The closed one-form
$A_s=X^\flat/W$ therefore has a global smooth primitive $\Phi$
on the open exterior.  Define the original slice by
\begin{equation}\label{eq:static:graph-embedding}
 \iota(p)=\bigl(t=-\Phi(p),\;r(p),\;\vartheta(p)\bigr).
\end{equation}
It is an injective graph over the entire static base.  Its
induced metric is
$h-Wd\Phi^2=h-WA_s^2=g$.  Moreover, the coordinate change
$t=z-\Phi$ transforms Equation~\eqref{eq:static:stationary-metric}
to the Schwarzschild--Tangherlini metric and takes the original
slice $z=0$ to Equation~\eqref{eq:static:graph-embedding}.
It carries its future normal to the future normal of this graph.
Consequently its second form is exactly
$-\sym\nabla X/u=K$.

At infinity the base metric is uniformly comparable to the
original end metric, and its Euclidean coordinates are obtained
by a fixed invertible linear map and a sublinear correction.
Thus $r\to\infty$ at the original end.  The graph stays
uniformly spacelike there: $u\to b^0<\infty$ and $W\to1$, while
\[
 W|d\Phi|_h^2=\frac{|X|_g^2}{u^2}=1-\frac W{u^2}.
\]
Consequently $|d\Phi|_h$ is bounded by a constant strictly less
than one sufficiently far out.  Integrating along base radial
curves gives $|t|\leq c_1r+C$ for some $c_1<1$.
Thus this end approaches the specified spatial infinity,
with an arbitrary admissible asymptotic boost permitted.

We verify smoothness at the original boundary; the two lapse
cases require different coordinates.  First, the static-base
isometry and its angular variables extend smoothly to the
one-sided attachment.  Away from the filled point the doubled
metric is $C^{2,\alpha}$, so its harmonic charts, and hence its
Euclidean isometry, are $C^{3,\alpha}$ in the collar coordinates.
On either closed side the metric coefficients are smooth.  The
isometry equation for its Cartesian components $Y^a$ reads
$\partial_i\partial_jY^a=\Gamma_{ij}^\ell\partial_\ell Y^a$;
differentiating this identity bootstraps the one-sided extension
to every order.  Normal geodesics for the static base metric $h$
from the seam then show that
its angular coordinates are the smooth boundary footpoint map.
In the positive boundary-lapse case the reflection in
$\lambda$ is an isometry of $h$ fixing the seam.  The angular
footpoint map is invariant under this reflection, hence is a
smooth even function of $\lambda$.  Taylor's formula for an
even smooth function shows it is smooth as a function of
$W=\lambda^2$ up to $W=0$.  Therefore the angular variables
are smooth in the original $(W,y)$ coordinates.  The areal
radius is also a smooth function of $W$ near zero, since
$W=1-2m/r^{n-2}$ has nonzero derivative at $r=r_h$.

Suppose $u>0$ at $S$.  The expression for $\alpha$ in the proof
of Proposition~\ref{prop:static:common-base} gives
\begin{equation}\label{eq:static:log-primitive}
 A_s=\frac1{2\kappa}\,d\log W+\gamma,
 \qquad \gamma\text{ smooth up to the original }S.
\end{equation}
Indeed $(b(W,y)-b(0,y))/W$ is smooth, and
$b(0,y)=1/(2\kappa)$ is constant.  The remainder $\gamma$
is closed.  Thus
$\Phi-(2\kappa)^{-1}\log W$ extends smoothly locally, and
these local extensions agree with the given global primitive
up to constants, so give a global smooth extension on a collar.
For the tortoise coordinate $dr_*/dr=W^{-1}$,
\begin{equation}\label{eq:static:tortoise}
 r_*=(2\kappa)^{-1}\log W+\gamma_0(W),
 \qquad \gamma_0\text{ smooth near }0.
\end{equation}
Here $W'(r_h)=(n-2)/r_h=2\kappa$; subtracting the displayed
simple pole from $dr_*/dW$ leaves a smooth coefficient.
The ingoing null coordinate
\[
 v=t+r_*=-\Phi+r_*
\]
therefore extends smoothly to $S$.  In $(v,r,\vartheta)$ the
spacetime metric is
$-Wdv^2+2dv\,dr+r^2\sigma_k$, smooth at $r_h$.
The boundary has finite ingoing time and $t\to+\infty$ as
$W\downarrow0$; it is a section of the future horizon.

Suppose instead $u=0$ on $S$.  Equation~\eqref{eq:static:zero-boundary}
shows that $A_s$ itself extends smoothly in the original collar,
and $\lambda$ is a smooth defining function there.  Consequently
$\Phi$ is smooth up to the boundary.  Equations~\eqref{eq:static:tortoise}
and $W=\lambda^2$ imply
$e^{\kappa r_*}=\lambda e^{\kappa\gamma_0(\lambda^2)}$.
The Kruskal coordinates
\begin{equation}\label{eq:static:kruskal}
 U=-e^{-\kappa(t-r_*)}
   =-\lambda e^{\kappa\Phi+\kappa\gamma_0(\lambda^2)},\qquad
 V=e^{\kappa(t+r_*)}
   =\lambda e^{-\kappa\Phi+\kappa\gamma_0(\lambda^2)}
\end{equation}
are smooth original functions and vanish together on $S$.
Their spacetime metric coefficient is a nonzero smooth function
of $UV$, by the simple zero of $W$ at $r_h$.
Thus this is a smooth extension to the bifurcation sphere,
with $U<0<V$ on the exterior side.  These signs agree with
the time orientation chosen in
Equation~\eqref{eq:static:stationary-metric}.

In either case the extended pullback metric equals $g$ by
continuity from the interior.  It is positive definite, so the
extended map is a spacelike immersion at every boundary point.
Its future unit normal is smooth in these regular spacetime
coordinates, and its second form extends as $K$ by the same
continuity.  The angular map on $S$ is a diffeomorphism onto
the full horizon sphere.  Distinct interior points were already
separated by their base coordinates; boundary points are
separated by their angles, and no interior point has radius
$r_h$.  Finally the original compact-core property and escape
of the end make this injective immersion proper.  It is
therefore a global smooth embedding with boundary, which
completes the proof of Theorem~\ref{thm:static:classification}.
